\section{为什么 Embodied Intelligence 比聊天更难}

虚拟世界里的错误通常表现为错误文字；物理世界里的错误会改变对象、损坏硬件，甚至伤害人。具身智能因此不仅需要识别“这是什么”，还要预测“我做了这个动作后会发生什么”。这把问题从静态语义理解扩展为闭环决策与后果建模。

\subsection{失败案例揭示的是 world model 缺口}

本节先用两个失败案例确定后续模型必须补上的能力。讲者展示机器人把可乐罐放入水槽时意外打开水龙头，以及机械臂收回时把容器倒置的例子。这些错误并非简单的 object detection failure；机器人可能看到了罐子、水槽和机械臂，却没有充分建模动作序列对液体、重力、碰撞和自身姿态的影响。

\lecturefigure{slide-02-why-embodied-intelligence.jpg}{真实世界的复杂性与 ambient intelligence 目标}{官方录像恢复幻灯片 V002；00:04:06}

可以用最小状态转移式表达这种缺口：
$$
s_{t+1}=f(s_t,a_t,\xi_t),\qquad o_t=h(s_t)+\epsilon_t.
$$
其中 $s_t$ 是时刻 $t$ 的真实环境状态，$a_t$ 是机器人动作，$f$ 是受动力学与接触影响的状态转移，$\xi_t$ 表示未建模扰动；$o_t$ 是传感器 observation，$h$ 是观测映射，$\epsilon_t$ 是测量噪声。机器人只看到 $o_t$，却必须推断隐藏的 $s_t$ 并预测动作后果。

\begin{knowledgebox}{First use：world model}
\term{World model（世界模型）}是对环境状态、动作与未来结果关系的内部表示。它可以是显式动力学模型、隐式 latent transition、视频预测器，也可以部分编码在 policy 中。它的价值不在于复刻整个世界，而在于支持与任务有关的预测、规划和错误检测。
\end{knowledgebox}

\teachervoice{Fei Xia 用“机器人把罐子放入水槽，却打开水龙头”的意外行为提醒学生：真实世界会产生训练者没有预想到的动作后果。这个例子不是可爱的 blooper，而是在说明 policy 若缺少 consequence model，随机误差也可能进入危险状态。}

\subsection{Gibson：智能来自 agent 与环境的闭环}

前面的失败自然引出 James J. Gibson 的 ecological view。与其只问模型内部存了什么，不如问 agent 被放进怎样的环境、能执行哪些动作、会收到什么反馈。儿童通过多年身体互动学习重力、摩擦、支撑、可达性和物体用途；机器人若没有对应经验，就难以仅靠少量 teleoperation trajectory 获得同样常识。

\lecturefigure{slide-03-gibson-ecological-perception.jpg}{Gibson 的生态知觉：理解 head 所处的环境}{官方录像恢复幻灯片 V003；00:05:18}

\begin{knowledgebox}{First use：affordance}
\term{Affordance（可供性）}描述环境相对于某个 agent 能做什么，例如把手“可抓”、平面“可放”、按钮“可按”。它不是物体的纯视觉属性，而取决于机器人形态、动作空间与任务。相同杯子对两指夹爪、人手和移动底盘具有不同 affordance。
\end{knowledgebox}

\subsection{Simulation 是经验放大器，不是真实世界替身}

安全 playpen 的工程实现之一是 simulation。讲者回顾自己在 iGibson 等工作中构建大规模室内场景、渲染 observation，并让 agent 在可交互环境中学习。Simulation 的优势是可重复、可并行、可访问隐藏状态并自动生成 label；它的风险是视觉、动力学、接触和任务分布与真实世界存在 gap。

\lecturefigure{slide-04-simulation-digital-twin-perception.jpg}{Simulation 同时提供 rendered observation、physics 与 semantic labels}{官方录像恢复幻灯片 V004；00:06:36}

\lecturefigure{slide-05-large-realistic-scenes.jpg}{大规模 realistic scenes：从几何资产走向可交互环境}{官方录像恢复幻灯片 V005；00:07:24}

\begin{warningbox}{Simulation 不自动等于 world model}
模拟器包含开发者写入的动力学与资产分布；agent 在模拟器里成功，可能只是利用了纹理、碰撞或 reset 机制的 shortcut。读 simulation result 时要区分：环境是否覆盖真实变化、传感器噪声是否合理、接触模型是否可信，以及 policy 是否在真实机器人上经过独立验证。
\end{warningbox}

\subsection{长时程任务暴露数据和组合性瓶颈}

单步抓取可以通过大量重复数据学习，长时程家务却要求在多个 room、object、failure recovery 与 subgoal 之间维持状态。图中的 household demo 不是为了证明家庭机器人已经解决，而是显示传统 imitation-learning pipeline 在场景规模、交互次数和任务组合上迅速变得昂贵。

\lecturefigure{slide-06-long-horizon-home-tasks.jpg}{长时程家庭任务：动作链、状态变化与失败恢复}{官方录像恢复幻灯片 V006；00:08:15}

\lecturefigure{slide-07-internet-ai-vs-embodied-ai.jpg}{Internet AI 与 Embodied AI：共享语义模块但动作闭环不同}{官方录像恢复幻灯片 V007；00:09:12}

\lecturefigure{slide-08-simulation-to-foundation-models.jpg}{从 50k episodes / 1.25M interactions 转向 Foundation Models}{官方录像恢复幻灯片 V008；00:11:15}

\teachervoice{讲者把自己从 large-scale simulation 转向 foundation models 的经历描述成一次研究路线变化：仅扩大场景与交互仍很难覆盖真实世界，而 web-scale 模型可能提供物体、语言和常识 prior。这里的关键词是“补充 prior”，不是“simulation 和 robot data 都不再需要”。}

\subsection{本章小结}

具身智能的难点来自部分可观测状态、连续动作、物理后果与长时程组合。Simulation 能放大经验，但仍有 sim-to-real gap；foundation models 的机会在于导入大规模视觉语言先验，下一章要回答这些先验怎样与低层控制相接。

